\subsection{Local systems over geometrically unibranch bases}

\begin{theorem}[C-RD7: effectivity over reduced geometrically unibranch bases]
\label{thm:C-RD7}
\status{Proved.}
Let $S$ be a reduced, locally noetherian, geometrically unibranch scheme.
In particular, $S$ may be any normal locally noetherian scheme.
Every system $A(s)$ of
abelian schemes over $\operatorname{Spec}\mathcal O_{S,s}$, equipped
with the compatible generization isomorphisms in C.4, is induced by an
abelian scheme over $S$, uniquely up to the unique isomorphism
respecting the entire system. No excellence hypothesis is required.
\end{theorem}
\begin{proof}
We first record the exact extension statement used here. If $X$ is a
reduced, geometrically unibranch, locally noetherian integral scheme
with generic point $\eta$,
and $A,B$ are abelian schemes over $X$, then
\[
 \operatorname{Hom}_X(A,B)\longrightarrow
 \operatorname{Hom}_\eta(A_\eta,B_\eta)
\]
is a bijection. Here is a proof which does not require finite
normalization. Work over a noetherian affine open in $X$. The scheme
$H=\operatorname{Hom}_{X\text{-}\mathrm{gp}}(A,B)$ is separated,
unramified, and essentially proper over $X$, meaning that it satisfies
the valuative criterion of properness
\cite[p.64]{Grothendieck1966Hom}. A generic homomorphism spreads over
some dense open $V\subset X$. In the proof of
\cite[Lemma~1.2.1, p.65]{Grothendieck1966Hom}, Grothendieck proves
explicitly that the reduced closure $Z$ of this section in $H$ is
finite and surjective over $X$. This cited closure assertion is
essential: essential properness alone is not being substituted for
properness of the possibly non-quasicompact scheme $H$.

The scheme $Z$ is integral and birational over $X$. On affine opens,
its finite coordinate algebra is a subalgebra of the function field,
and therefore of the integral closure of $\mathcal O_X$. Consequently
the normalization of $X$ maps surjectively to $Z$. Geometric
unibranchness says that this normalization map is a universal
homeomorphism. It follows, by surjectivity after every field extension,
that $Z\to X$ is radicial. As a closed subscheme of $H$, it is also
unramified. A finite radicial unramified morphism is a closed immersion;
its schematic dominance, or simply reducedness and birationality of
$X$, makes that immersion an isomorphism. Thus the section extends.
Separatedness gives uniqueness. Applying the result to a generic
isomorphism and its inverse extends generic isomorphisms as well.

The irreducible components of a locally noetherian geometrically
unibranch scheme are locally disjoint and open. Work on one such
component, which is integral because $S$ is reduced.
Let $\eta$ be its generic point. Since an abelian scheme is of finite
presentation, each $A(s)$ spreads to an abelian scheme $B_s$ over
some open neighbourhood $V_s$ of $s$; the group operations, properness,
smoothness and geometrically connected fibres spread after shrinking.
The given generization maps identify every $(B_s)_\eta$ with the
fixed abelian variety $A(\eta)$. On $V_s\cap V_t$, the resulting
generic isomorphism extends uniquely by the preceding paragraph.
Uniqueness proves the cocycle condition. These isomorphisms glue the
$B_s$ to an abelian scheme $A/S$.

For a point $t$, compare $A$ over $\operatorname{Spec}\mathcal O_{S,t}$
with the prescribed $A(t)$. Both have the stipulated generic
identification with $A(\eta)$; the same Hom extension theorem
therefore gives the unique local isomorphism. Its compatibility with
all generization maps follows from equality at the generic point.
Finally these constructions glue over the locally disjoint components
of $S$.
\end{proof}

\begin{lemma}[Compatible local sections of an unramified scheme]
\label{lem:C-RD8-sections}
Let $X$ be locally noetherian and let $H\to X$ be separated,
unramified and locally of finite presentation. Suppose a section over
$\operatorname{Spec}\mathcal O_{X,x}$ is given for every $x\in X$,
compatibly with all generizations. Then these sections are induced by
a unique section over $X$.
\end{lemma}
\begin{proof}
Work on a noetherian open of $X$. Spread the section at $x$ to an open
neighbourhood $V_x$. Delete from $V_x$ all irreducible components of
$X$ which do not contain $x$. Every generic point of the remaining
$V_x$ then generalizes $x$, so the spread has the prescribed generic
value. On an overlap, two spreads agree at every generic point.
Separatedness makes them agree on the reduction of the overlap,
because that reduction is schematically dominated by its generic
points. The nilradical is nilpotent locally; formal unramifiedness
makes the two spreads agree on the entire overlap. They glue, and
the same argument, or localization at every point, gives uniqueness.
\end{proof}

\begin{theorem}[C-RD8: reduced bases with a finite geometrically unibranch cover]
\label{thm:C-RD8}
\status{Proved, with the covering hypothesis displayed explicitly.}
Let $S$ be a reduced locally noetherian scheme which, locally on $S$,
admits a finite surjective morphism $S'\to S$ with $S'$ geometrically
unibranch. Then every local system in C.4 is effective, uniquely up to
the unique isomorphism respecting its prescribed identifications.
The normalization of $S$ need not be finite.
\end{theorem}
\begin{proof}
The assertion is local on $S$, with unique compatible gluing, so fix
such a cover. Replace $S'$ by its reduction. Pull the local system
back by assigning to $x\in S'$ the base change of $A(f(x))$ to
$\operatorname{Spec}\mathcal O_{S',x}$. Theorem~\ref{thm:C-RD7}
produces an abelian scheme $B'/S'$ with all these local identifications.

On $D=S'\times_SS'$, these identifications give, at every point $z$
lying over $s\in S$, an isomorphism between the two pullbacks of $B'$:
both are identified with the base change of $A(s)$. These local
isomorphisms are generization-compatible. Their Isom scheme is
separated, unramified and locally of finite presentation, so
Lemma~\ref{lem:C-RD8-sections} gives a global isomorphism on $D$.
The same lemma on the triple fibre product verifies its cocycle
condition. This argument applies even when these fibre products are
nonreduced.

Fix $s\in S$. On $S'\times_S\operatorname{Spec}\mathcal O_{S,s}$,
compare $B'$ with the pullback of $A(s)$. At every local point the
given generization maps and the local identifications of $B'$ give
compatible isomorphisms. The lemma globalizes them on this scheme.
Spread $A(s)$ to an abelian scheme $B_s$ on a neighbourhood $V_s$.
Finite presentation of the schemes and their isomorphism, together
with finiteness of $S'\to S$, spreads the last isomorphism to
\[
 B_s\times_S S'\simeq B'|_{f^{-1}(V_s)}
\]
after shrinking $V_s$. This step spreads one isomorphism of abelian
schemes of finite presentation, and does not spread an infinite BT
system. Over $D\times_S\operatorname{Spec}\mathcal O_{S,s}$ this
isomorphism respects the descent isomorphism just constructed. Since
$D\to S$ is finite, that equality of finitely presented morphisms
also spreads after a further shrinking of $V_s$. Thus the chosen
isomorphism over $f^{-1}(V_s)$ respects the descent isomorphism
throughout $D|_{V_s}$.

It remains to descend the resulting isomorphisms between $B_s$ and
$B_t$ on overlaps. They satisfy the descent equality on $D$, by their
pointwise construction and the preceding lemma. The finite surjection
$S'\to S$ is schematically dominant because $S$ is reduced. Apply
Proposition~\ref{prop:C-RD1} to their BT homomorphisms, obtaining a
unique BT homomorphism over $V_s\cap V_t$. At every generic point
this map is algebraic: faithfully flat descent over the finite
nonzero algebra above that field descends the abelian homomorphism.
Every associated point of the reduced overlap is generic, so
\cite[Proposition~1.2, p.64]{Grothendieck1966Hom} makes the descended
BT map algebraic on the whole overlap. Apply the same argument to its
inverse. Faithfulness verifies both inverse identities and the
cocycle. The $B_s$ therefore glue to an abelian scheme over $S$.

Its local identification with $A(s)$ is already built into $B_s$.
All prescribed generization compatibilities and uniqueness follow
either after the schematically dominant cover or from faithfulness
of the BT functor. This proves the theorem.
\end{proof}

\begin{remark}[The nodal involution does not supply the advertised counterexample]
The proposed gluing by $[-1]$ at the two branches of a node has finite
monodromy and is effective: twist the constant abelian scheme by the
corresponding finite etale double cover. This observation refutes that
specific proposed construction, but does not by itself settle C.4 on
all reduced or all nonreduced noetherian bases. In particular, a proof
which shrinks to a normal dense open must justify that such an open
exists; this cannot be silently assumed for an arbitrary nonexcellent
noetherian integral scheme.
\end{remark}
